\documentclass[12pt,a4paper,oneside]{report}
\usepackage{indentfirst}
\usepackage{times}
\setlength\parindent{1cm}
\renewcommand{\baselinestretch}{1.50}\normalsize
\usepackage{anysize}

\marginsize{1.25in}{.50in}{0.25in}{1in}
\usepackage{cite}
\usepackage[colorlinks=true, linkcolor=black, citecolor=black, urlcolor=black]{hyperref}

\usepackage{graphics}
\usepackage{graphicx}
\usepackage{epsfig}
\usepackage[fleqn]{amsmath}
\usepackage{amsfonts}
\usepackage[section]{placeins}
\usepackage{algorithm}
\usepackage{algpseudocode}
\usepackage{pifont}
\usepackage[utf8]{inputenc}
\usepackage[T1]{fontenc}
\usepackage[force,almostfull]{textcomp}
\usepackage{graphicx}
\usepackage{setspace}
\usepackage{fancyhdr}
\usepackage{truncate}
\usepackage{nomencl} 
\usepackage{acronym}
\usepackage{array}
\usepackage{caption}
\usepackage{subcaption}
\usepackage[overload]{textcase}
\usepackage{listings}
\usepackage[paperwidth=612pt, paperheight=792pt,top=72pt, right=72pt,bottom=72pt,left= 72pt]{geometry}
\usepackage{booktabs}
\usepackage{adjustbox}
\usepackage{enumitem}
\usepackage{longtable}
\usepackage{multicol}

\renewcommand{\nomname}{List of Abbreviations}
\usepackage{makeidx}
\makeindex
\makenomenclature
\newcommand{\quotes}[1]{``#1''}
\usepackage{titlesec}
\titleformat{\chapter}[display]
{\normalfont\Large\bfseries\centering}
{\chaptertitlename\ \thechapter}{15pt}{\LARGE}
\titleformat{\section}{\large\bfseries}{\thesection}{1em}{}
\titleformat{\subsection}{\normalsize\bfseries}{\thesubsection}{1em}{}
\renewcommand{\chaptermark}[1]{\markboth{ \emph{#1}}{}}

\printnomenclature[5em]
\pagestyle{fancy}
\fancyhf{}% Clears all header and footer fields
\renewcommand{\footrulewidth}{2pt}
\renewcommand{\headrulewidth}{1.2pt}
\rhead{\scriptsize {\leftmark}}
\lhead{\small{College of Engineering, Cherthala  \;\;\;\;\;\;\;\;\;\;\;\;\;\;\;\;\;\;}}
\rfoot{\thepage}
\cfoot{\empty}
\lfoot{\small{Dept of Computer Sci \& Engg  \;\;\;\;\;\;\;\;\;\;\;\;\;\;\;\;\;\;}}
%\fancyfoot[L]{My Document Footer}
\renewcommand{\figurename}{Fig.}
\renewcommand\bibname{References}
\setlength{\arrayrulewidth}{1mm}
\setlength{\tabcolsep}{18pt}
\renewcommand{\arraystretch}{1.5}

\begin{document}
\thispagestyle{fancy}
\pagestyle{fancy}

\begin{titlepage}
\begin{center}
\Large{\textbf{MAIN PROJECT REPORT ON}}

\begin{singlespace}
\LARGE{\textbf{STYLMORPH: AI-POWERED VIRTUAL TRY-ON PLATFORM}}\\
\end{singlespace}
\Large{\textit{Submitted By }}\\
\large{\textit{\textbf{Devi Nandana}
(\textbf{Reg. No. CEC22CS054})}}\\
\large{\textit{\textbf{Neeraj S}
(\textbf{Reg. No. CEC22CS100})}}\\
\large{\textit{\textbf{Sanjay Kumar P M}
(\textbf{Reg. No. CEC22CS114})}}\\
\large{\textit{\textbf{Vinayakrishnan U}
(\textbf{Reg. No. CEC22CS128})}}\\

\large{\textit{under the esteemed guidance of}}\\
\large{\textit{\textbf{Ms. Jincy Showri} }}\\
\large{\textit{Assistant Professor}}\\
\large{\textit{Department of Computer Science and Engineering}}
\\
\begin{figure}[h]
\begin{center}

\epsfig{width=1.2 in, file=logo.jpg}
\end{center}
\end{figure}
\begin{singlespace}
\large{\textbf{MARCH 2026}}\\

\large{\textbf{DEPARTMENT OF COMPUTER SCIENCE AND ENGINEERING\\COLLEGE OF ENGINEERING, CHERTHALA\\ PALLIPPURAM P O, ALAPPUZHA-688541, \\PHONE: 0478 2553416, FAX: 0478 2552714\\http://www.cectl.ac.in}}
\end{singlespace}
\end{center}
\end{titlepage}

\begin{titlepage}
\begin{center}
\large\textbf{MAIN PROJECT REPORT ON}\\
\begin{center}
\textbf{\Large{STYLMORPH: AI-POWERED VIRTUAL TRY-ON PLATFORM}}
\end{center}

\large{\textit{Submitted By }}\\
\begin{singlespace}
\large{\textit{\textbf{Devi Nandana}
(\textbf{Reg.No CEC22CS054})}}\\
\large{\textit{\textbf{Neeraj S}
(\textbf{Reg. No. CEC22CS100})}}\\
\large{\textit{\textbf{Sanjay Kumar P M}
(\textbf{Reg. No. CEC22CS114})}}\\
\large{\textit{\textbf{Vinaya Krishnan U}
(\textbf{Reg. No. CEC22CS128})}}\\
\end{singlespace}
\large{\textit{under the esteemed guidance of}}\\
\large{\textit{\textbf{Ms. Jincy Showri}}}\\
\begin{singlespace}
\large{\textit{In partial fulfillment of the requirements for the award of the degree}\\
\large{ \textit{of}}\\
\large{\textit{Bachelor of Technology} }\\
\large{\textit{in}}\\
\large{\textit{Computer Science and Engineering}}\\
\large{\textit{of}}\\
\large{\textit{A P J Abdul Kalam Technological University }}}\\
\end{singlespace}
\begin{figure}[h]
\begin{center}

\epsfig{width=1.2in, file=logo.jpg}
\end{center}
\end{figure}
\vspace{-.2in}
\begin{singlespace}
\large{\textbf{MARCH 2026\\Department of Computer Science And 
Engineering\\College of Engineering, Pallippuram P O, Cherthala\\ Alappuzha Pin: 688541\\
Phone: 0478 2553416, Fax: 0478 2552714 \\ http://www.cectl.ac.in}}
\end{singlespace}
\end{center}
\end{titlepage}

\begin{titlepage}
\begin{center}
\large{\textbf{DEPARTMENT OF COMPUTER SCIENCE AND ENGINEERING}}\\
\large{\textbf{COLLEGE OF ENGINEERING CHERTHALA\\ALAPPUZHA-688541}}\\
\end{center}
\begin{figure}[h]
\begin{center}

\epsfig{width=1.5in, file=logo.jpg}
\end{center}
\end{figure}
\begin{center}
\large{\textbf{CERTIFICATE}}\\
\end{center}
\begin{spacing}{1.5}
This is to certify that, the main project report on \textbf{\textit{STYLMORPH: AI-POWERED VIRTUAL TRY-ON PLATFORM}} is a bonafide record of the \textbf{CSD415 Major Project Phase II} presented by \textbf{Devi Nandana} (CEC22CS054), \textbf{Neeraj S} (CEC22CS100), \textbf{Sanjay Kumar P M} (CEC22CS114), \textbf{Vinaya Krishnan U} (CEC22CS128), Seventh Semester B. Tech. Computer Science and Engineering students, under our guidance and supervision, in partial fulfillment of the requirements for the award of the degree, \textbf{B. Tech. Computer Science and Engineering} of \textbf{A P J Abdul Kalam Technological University}.
\end{spacing}
\begin{tabbing}
xxxxxxxxxxxxxxxxxxxxxxxxxxxxxxxxxxxxxx\= xxxxxxxxxxxxxxxxxxxx\= \kill
\hspace{.15in}{\bf Guide}\>\hspace{-.70in}{\bf{Co-ordinator}} \>\hspace{-.16in}{\bf HoD} \\
\\
\end{tabbing}
\begin{tabbing}
xxxxxxxxxxxxxxxxxxxxxxxxxxxxxxxxxxxxxx\= xxxxxxxxxxxxxxxxxxxxx\= \kill
\hspace{.15in}{\bf{Ms. Jincy Showri}}\>\hspace{-.75in} {\bf{Dr. Vishnupriya G S}} \> \hspace{-.24in}{\bf{Dr. Preetha Theresa Joy}}\\
\hspace{.1in} Assistant Professor\>\hspace{-.70in}{{Assistant Professor}}\>\hspace{-.24in}Professor\\\hspace{.1in} Dept. of Computer Sci \& Engg \>\hspace{-0.70in}Dept. of Computer Sci \& Engg\>\hspace{-.28in} Dept. of Computer Sci \& Engg\\
\end{tabbing}
\end{titlepage}

\renewcommand*\thesection{\thechapter.\arabic{section}}
\begin{titlepage}
\pagenumbering{roman}
\setcounter{page}{4}
\section*{\begin{center} \Large{ACKNOWLEDGEMENT}\end{center}}
\thispagestyle{plain}

\par
This work would not have been possible without the support of many people. First and the foremost, we thank Almighty God who gave us the inner strength, resource and ability to complete our project successfully.

\par
We would like to thank \textbf{Dr. Jaya V L}, The Principal, who has provided with the best facilities and atmosphere for the project completion and presentation. We would also like to thank our HoD \textbf{Dr. Preetha Theresa Joy} (Professor, Computer Engineering), Project Coordinator \textbf{Dr. Vishnupriya G S} (Assistant Professor, Computer Engineering) and Project Guide \textbf{Ms. Jincy Showri} (Assistant Professor, Computer Engineering) for the help extended and also for the encouragement and support given to us while doing the project.

\par
We would like to thank our dear friends for extending their cooperation and encouragement throughout the project work, without which we would never have completed the project this well. Thank you all for your love and also for being very understanding.
\end{titlepage}

\newpage
\pagestyle{plain}

\begin{center}
    \Large
    \textbf{ABSTRACT}
\end{center}

\vspace{0.5cm}

\noindent
STYLMORPH is an innovative AI-powered virtual try-on platform that revolutionizes online fashion retail by enabling users to visualize garments on their own photographs and generate personalized 3D avatars from single 2D images. The system integrates a robust two-stage processing workflow combining high-speed 2D virtual try-on using the FASHN-VTON framework with detailed 3D avatar generation using PIFuHD monocular 3D reconstruction. By providing both immediate 2D garment previews and detailed 3D avatars for comprehensive 360-degree fitting analysis, Stylmorph significantly mitigates size and fit discrepancies. The MERN-stack architecture, augmented by Three.js WebGL rendering and Cloudinary CDN delivery, proves that high-fidelity fashion visualization is viable for real-world e-commerce deployment. This interactive and personalized shopping experience aims to reduce the high product return rates burdening the retail sector and substantially augment consumer confidence in online apparel purchases.

\vspace{0.3cm}

\noindent
\textbf{Keywords:} Virtual Try-On, 3D Avatars, Monocular 3D Reconstruction, PIFuHD, FASHN-VTON, E-commerce, Computer Vision, Deep Learning, Three.js

\tableofcontents
\listoffigures
\listoftables

\newpage
\pagestyle{fancy}
\pagenumbering{arabic}
\chapter{INTRODUCTION}

The rapid growth of e-commerce has fundamentally transformed how consumers shop for clothing, yet significant challenges persist in providing accurate fit visualization and personalized shopping experiences online. At least 30-40\% of online fashion purchases are returned due to improper fit and inability to visualize clothing on individual body types. Traditional e-commerce platforms rely primarily on static 2D images and generic size charts, leaving customers uncertain about how garments will fit their unique body proportions and personal style preferences. This uncertainty contributes to purchase hesitation, high return rates, and overall customer dissatisfaction.

In the rapidly evolving landscape of e-commerce and digital fashion, STYLMORPH emerges as a comprehensive virtual try-on platform designed to address the lack of personalization and fit visualization in modern e-commerce. Early attempts at virtual fitting rooms relied on users manually entering extensive body measurements to generate generic parametric bodies. While functional, these systems lacked visual personalization and real-world fidelity. To address these limitations, STYLMORPH introduces a two-stage approach that combines 2D visualization with detailed 3D avatar generation.

The system operates through a sophisticated pipeline that begins with user input collection. Users can upload a single full-body image, and the system first applies 2D virtual try-on using FASHN-VTON to synthesize selected garments on the user's photograph, preserving garment details, textures, and folds while adapting to the target person's pose and body shape. Subsequently, the same image is used to generate a personalized 3D avatar using PIFuHD, a multi-level pixel-aligned implicit function that enables high-resolution 3D human digitization from a single image, preserving details such as clothing folds and hair.

STYLMORPH's architecture comprises multiple interconnected modules built on the MERN stack (MongoDB, Express.js, React.js, Node.js) with specialized Python microservices for machine learning processing. The platform combines an intuitive React.js frontend with an AI-driven backend. Due to infrastructure constraints, the FASHN-VTON 2D try-on pipeline is deployed on Google Colaboratory with a GPU runtime, accessed via secure API tunneling from the Node.js backend. The PIFuHD pipeline runs locally within a Python subprocess. Generated meshes are delivered rapidly via Cloudinary and rendered in the browser using Three.js, enabling hardware-accelerated, real-time interactive visualization of 3D avatars.

This comprehensive approach not only enhances user confidence in purchasing decisions but also reduces the environmental impact of modern retail by minimizing product returns and associated logistics. By accurately simulating how clothing items appear on individual body types, STYLMORPH significantly enhances customer confidence in online purchases, reduces return rates, and creates a more engaging shopping experience.

\chapter{PROBLEM STATEMENT}

The online fashion industry faces significant challenges with high return rates primarily due to improper fit and inability to visualize clothing on individual body types. Traditional e-commerce platforms provide limited tools for customers to assess how garments will look on their specific body shape, leading to purchase uncertainty and dissatisfaction. 

Current virtual try-on solutions often rely on generic parametric bodies like the SMPL model, which typically produces naked or minimally-clothed bodies stripped of hair and complex textures. These standardized representations fail to capture the nuanced details that consumers need to make confident purchasing decisions.

\section{Objective}

The key objectives of Stylmorph are:
\begin{itemize}
    \item To provide instant 2D try-on results for rapid garment browsing using FASHN-VTON
    \item To generate detailed, topologically accurate 3D avatars from a single photo without specialized hardware using PIFuHD
    \item To enable real-time, 360-degree interactive inspection of how garments fit the user's body within a standard web browser
    \item To create a scalable web architecture with MERN stack for seamless user experience
    \item To reduce return rates through accurate fit prediction and realistic visualization
\end{itemize}

\chapter{EXISTING SYSTEM}

Current virtual try-on solutions in the e-commerce landscape suffer from several limitations that restrict their effectiveness and adoption. Most existing systems rely on standardized body models that fail to account for individual variations in body shape, proportions, and measurements. These generic avatars provide inaccurate representations of how clothing will actually fit and appear on specific users.

\section{Current Virtual Try-On Solutions}

Existing virtual try-on platforms typically use one-size-fits-all approaches that cannot accurately simulate clothing fit across diverse body types. Research on virtual try-on systems has evolved significantly over the past decade, but early solutions primarily used 2D image overlays and basic body measurement calculators. While these systems provided basic functionality, they struggled with realistic volume simulation and true personalized fit accuracy.

The introduction of sophisticated 3D modeling and monocular reconstruction techniques marked a massive leap forward. Loper et al. introduced the Skinned Multi-Person Linear (SMPL) model, a vertex-based model that represents a wide variety of human body shapes accurately. However, SMPL typically produces naked or minimally-clothed bodies stripped of hair and complex textures, limiting its utility for realistic virtual try-on.

\section{Disadvantages of Existing System}

Current systems face the following challenges:
\begin{itemize}
    \item \textbf{Generic Avatar Models:} Use standardized body types that don't reflect individual user measurements, failing to capture detailed features like clothing wrinkles and hairstyles
    \item \textbf{Limited Customization:} Offer minimal options for body shape adjustments and rely on parametric constraints
    \item \textbf{Inaccurate Fit Prediction:} Cannot accurately simulate how different fabrics and cuts will drape on various body types
    \item \textbf{Poor Integration:} Limited connectivity with e-commerce platforms and often remain isolated research prototypes lacking seamless web deployment
    \item \textbf{Basic Visualization:} Provide 2D or simplified 3D representations that lack realism
    \item \textbf{Single-Modality Approaches:} Most systems offer either 2D try-on or 3D avatars, but rarely both in an integrated workflow
\end{itemize}

\chapter{LITERATURE SURVEY}

The field of 3D avatar generation and virtual try-on systems has witnessed significant advancement in recent years, driven by breakthroughs in computer vision, deep learning, and 3D reconstruction techniques. This chapter presents a comprehensive review of existing literature and research works relevant to the development of STYLMORPH, an AI-powered virtual try-on platform. The survey covers ten significant contributions that form the foundation for our work, ranging from parametric body models to neural rendering techniques.

\section*{1. PIFuHD: Multi-Level Pixel-Aligned Implicit Function for High-Resolution 3D Human Digitization (2020)}

\textbf{Authors: Shunsuke Saito, Zeng Huang, Ryota Natsume, Shigeo Morishima, Angjoo Kanazawa, Michael J. Black}

PIFuHD introduces a multi-level pixel-aligned implicit function that dramatically improved the resolution of 3D human digitization from a single image. The architecture employs a coarse-level network that observes the whole image for global context, and a fine-level network that focuses on local features to predict detailed surface geometry. This hierarchical approach enables the reconstruction of intricate details such as clothing wrinkles and hairstyles, generating realistic meshes without parametric constraints. The primary advantage of this method is its ability to produce high-resolution output with fine detail preservation, capturing elements like wrinkles and hair without being constrained by parametric models. It also works with a single image input, eliminating the need for multi-camera setups. However, the approach is computationally intensive, requiring significant processing time and a high-end GPU for inference. The quality of the reconstruction relies heavily on the input image's lighting, pose, and occlusions, and the model has a large memory footprint for high-resolution outputs. Furthermore, its complex training pipeline requires substantial data.

\section*{2. FASHN-VTON: High-Fidelity Virtual Try-On with Texture Preservation (2024)}

\textbf{Authors: FASHN Research Team}

FASHN-VTON represents a novel approach for high-fidelity virtual clothing transfer that preserves garment details, textures, and folds while adapting to the target person's pose and body shape. This framework builds upon earlier work in 2D virtual try-on such as VITON and CP-VTON, establishing a robust foundation for image-based clothing transfer with superior texture preservation and garment morphing capabilities. Its key advantages are high-fidelity texture preservation that maintains garment logos and patterns, accurate garment morphing that adapts to the target person's pose, and its ability to adapt to varied body shapes and poses effectively. It also preserves fine details like folds and fabric textures. On the downside, the method is computationally intensive, requiring GPU acceleration, and remote execution can introduce significant latency of 7-8 minutes. It requires high-quality input images for optimal results and is limited to 2D visualization without any depth information, with a dependency on external compute resources.

\section*{3. AG3D: Learning to Generate 3D Avatars from 2D Images (2023)}

\textbf{Authors: Zijian Dong, Xu Chen, Jinlong Yang, Michael J. Black, Otmar Hilliges, Andreas Geiger}

AG3D introduces a method for generating detailed 3D avatars from 2D images via Fast-SNARF, utilizing tri-plane features for detailed geometry reconstruction. The approach employs adversarial training to enhance realism and enables the generation of avatars without requiring 3D scans for training. Its main advantages are that it requires no 3D scans for training, working directly with 2D image datasets, while producing realistic avatars with detailed surface geometry and capturing clothing details effectively. Fast-SNARF enables efficient deformation handling, and tri-plane features provide an efficient geometry representation. However, it needs a high-end GPU for inference and is not suitable for real-time applications. The complex training pipeline requires expertise, and the results are limited by the quality of the input images. Additionally, the adversarial training process can be unstable.

\section*{4. CAPE: Clothed Auto-Personalized Avatars (2022)}

\textbf{Authors: Qianli Ma, Jinlong Yang, Anurag Ranjan, Sergi Pujades, Gerard Pons-Moll, Michael J. Black, Siyu Tang}

CAPE extends the SMPL model with a generative clothing layer, using a VAE-GAN architecture for realistic cloth deformation. The model provides pose-dependent garment adaptability, enabling realistic wrinkles and folds that respond to body movements. The system excels at generating realistic wrinkles and folds and works seamlessly with the standard SMPL model. It adapts to body poses with dynamic cloth deformation, and its generative approach enables varied clothing styles with pose-dependent modeling that enhances realism. The main disadvantages include its requirement for high-quality 3D scans for training and its limitation to garment types seen during training. Data collection and processing are expensive, and the inference is computationally intensive. Furthermore, it requires careful balancing of GAN training.

\section*{5. SMPL: A Skinned Multi-Person Linear Model (2015)}

\textbf{Authors: Matthew Loper, Naureen Mahmood, Javier Romero, Gerard Pons-Moll, Michael J. Black}

SMPL introduces a parametric human body model for animation, utilizing linear blend skinning and shape parameters to provide realistic body deformations. The model has become a widely adopted standard in computer vision and graphics research, serving as the foundation for numerous subsequent works. Its advantages include being a highly efficient parametric representation that integrates easily with animation pipelines. It has seen wide adoption in the research community, offers realistic body deformations with pose changes, and provides a compact representation through shape and pose parameters. The disadvantages are that it has limited clothing and hair detail, typically producing naked or minimally-clothed bodies. It requires parameter tuning for accuracy and lacks fine surface details and textures, making it unable to capture clothing wrinkles or hairstyles.

\section*{6. PIXIE: Portrait Image to Full-Body Avatar (2021)}

\textbf{Authors: Timo Bolkart, Jinlong Yang, Michael J. Black}

PIXIE generates full-body 3D models from single portrait images using deep shape and pose priors, combining facial and body reconstruction networks. The approach enables fast inference while capturing realistic body shapes from minimal input. It requires only a single image input, has a fast inference time suitable for applications, captures realistic body shape variations, combines facial and body reconstruction, and uses deep priors to ensure plausible shape generation. However, it is less accurate for occluded body parts and offers limited geometric detail for clothing. It is trained on specific dataset distributions and can struggle with extreme poses, and it may not capture fine surface details.

\section*{7. MetaHuman Creator: Ultra-Realistic Digital Humans (2021)}

\textbf{Authors: Epic Games Research Team}

MetaHuman Creator is a cloud-based system for creating ultra-realistic digital humans, offering real-time customization and facial rigging with direct integration into Unreal Engine. The platform represents the state-of-the-art in high-fidelity digital human creation. It achieves photorealistic human quality, offers easy customization through an intuitive interface, and provides real-time facial rigging and animation. It also offers cloud-based storage and processing with direct integration into game engines. The disadvantages are its heavy computational cost for rendering, its requirement of the Unreal Engine ecosystem, and its limited offline usage capabilities. It is not designed for single-photo reconstruction and is considered overkill for e-commerce try-on applications.

\section*{8. HumanNeRF: Neural Rendering for Dynamic Humans (2021)}

\textbf{Authors: Chung-Yi Wang, Tomas Simon, et al.}

HumanNeRF employs neural radiance fields for human reconstruction, handling motion with pose-conditioned NeRFs to produce high-fidelity dynamic humans. The approach enables photorealistic novel view synthesis of moving subjects. The method provides high visual realism with neural rendering, handles dynamic motion with pose-conditioning, produces high-fidelity dynamic human models, enables photorealistic novel view synthesis, and captures fine details like clothing folds. On the other hand, it has a slow rendering speed unsuitable for real-time, requires a complex neural network setup, is memory intensive for storage and processing, and its training requires multi-view video, as it is not designed for single-image inference.

\section*{9. Deep Fashion3D: Dataset for 3D Garment Reconstruction (2021)}

\textbf{Authors: Heming Zhu, Yu Cao, et al.}

Deep Fashion3D provides a large-scale dataset and benchmark for 3D garment reconstruction, including topology-consistent meshes specifically focused on fashion and virtual try-on applications. The dataset enables research in garment transfer and texture learning. It offers a large-scale dataset for clothing research with a topology-consistent mesh representation. It is specifically designed for fashion applications, supports texture learning and garment transfer, and serves as a benchmark for 3D garment reconstruction tasks. Its limitations include limited pose and garment variety, a requirement for large storage capacity, restricted license terms for commercial use, static poses that limit dynamic applications, and varying annotation quality across samples.

\section*{10. AvatarCLIP: Zero-Shot Text-Guided Avatar Generation (2022)}

\textbf{Authors: Fangzhou Hong, Mingyuan Zhang, et al.}

AvatarCLIP uses CLIP to align text with avatar shapes, generating avatars based on textual prompts through zero-shot learning for open vocabulary. The approach enables creative avatar design without dataset-specific training. Its advantages are that it requires no dataset-specific training, its text-guided generation enables creative design, its zero-shot learning works with an open vocabulary, it supports novel avatar concepts, and it is flexible and adaptable to user prompts. The disadvantages are that it offers limited fine control over avatar details, has lower geometric accuracy than specialized models, its output quality depends on the CLIP model, it may produce inconsistent results, and it cannot capture a specific user's likeness.

\section*{11. DreamHuman: Diffusion-Based Human Modeling (2023)}

\textbf{Authors: Jie Hu, et al.}

DreamHuman employs diffusion models for 3D human synthesis, combining multi-view diffusion and NeRFs to generate diverse realistic humans. The approach demonstrates strong generalization capabilities by integrating multiple state-of-the-art AI models. The strengths of DreamHuman are its high realism and diversity in generated humans, strong generalization across different human types, effective combination of multi-view diffusion with NeRFs, generation of diverse realistic human models, and successful integration of multiple AI models. The drawbacks are its very high computational cost, slow inference speed that limits applications, complex implementation requiring expertise, lack of optimization for single-image reconstruction, and the significant resources needed for training.






\section{Summary and Comparative Analysis}

Table 4.2 provides a comprehensive comparative summary of the surveyed literature, highlighting key contributions, methodologies, advantages, and limitations relevant to STYLMORPH's objectives.

\begin{center}
\vspace{-\baselineskip} % Reduce vertical space
\hspace{-0.7cm} % Shift table to the left
\small % Reduce font size to fit table better
\begin{longtable}{|p{0.1cm}|p{1.9cm}|p{2.0cm}|p{2.0cm}|p{2.0cm}|p{2.0cm}|}
\caption{Comprehensive Literature Survey Summary}
\label{tab:literature-summary}\\
\hline
\textbf{S.No.} & \textbf{Reference} & \textbf{Key Contribution} & \textbf{Methodology} & \textbf{Advantages} & \textbf{Disadvantages} \\
\hline
\endfirsthead

\hline
\textbf{S.No.} & \textbf{Reference} & \textbf{Key Contribution} & \textbf{Methodology} & \textbf{Advantages} & \textbf{Disadvantages} \\
\hline
\endhead

\hline
\multicolumn{6}{r}{\small{Continued on next page}} \\
\endfoot

\hline
\endlastfoot

1 & PIFuHD (Saito et al., 2020) & High-resolution 3D human digitization from single images & Multi-level pixel-aligned implicit functions & High-resolution output, preserves fine details, no parametric constraints & Computationally intensive, relies on image quality, requires significant processing time \\
\hline
2 & FASHN-VTON (FASHN Research, 2024) & High-fidelity 2D virtual try-on with texture preservation & Deep learning garment transfer with warping and synthesis & High-fidelity texture preservation, accurate garment morphing, pose adaptation & GPU acceleration required, remote execution latency (7-8 min) \\
\hline
3 & AG3D (Dong et al., 2023) & 3D avatar generation from 2D images & Fast-SNARF, tri-plane features, adversarial training & No 3D scans needed, realistic avatars, captures clothing details & High-end GPU needed, not real-time, complex training \\
\hline
4 & CAPE (Ma et al., 2022) & Clothed avatars with generative clothing layer & VAE-GAN for cloth deformation, SMPL extension & Realistic wrinkles, SMPL compatible, pose adaptation & Requires 3D scans, limited to trained garments, expensive training \\
\hline
5 & SMPL (Loper et al., 2015) & Parametric human body model & Linear blend skinning, shape and pose parameters & Highly efficient, easy animation integration, widely adopted & Limited clothing detail, needs parameter tuning, lacks fine surface realism \\
\hline
6 & PIXIE (Bolkart et al., 2021) & Full-body avatars from single portraits & Deep shape/pose priors, facial+body networks & Single image input, fast inference, realistic body shapes & Occlusion sensitivity, limited clothing detail, dataset-specific \\
\hline
7 & MetaHuman (Epic Games, 2021) & Ultra-realistic digital humans & Cloud-based system, real-time customization & Photorealistic quality, easy customization, real-time rendering & Heavy compute cost, Unreal Engine dependent, limited offline use \\
\hline
8 & HumanNeRF (Wang et al., 2021) & Neural rendering for dynamic humans & Pose-conditioned NeRFs, neural radiance fields & High visual realism, dynamic motion, photorealistic & Slow rendering, complex setup, memory intensive \\
\hline
9 & Deep Fashion3D (Zhu et al., 2021) & 3D garment dataset and benchmark & Topology-consistent meshes, fashion focus & Large clothing dataset, texture learning support, enables garment transfer & Limited pose variety, large storage, restricted license \\
\hline
10 & AvatarCLIP (Hong et al., 2022) & Zero-shot text-to-avatar generation & CLIP alignment, open vocabulary learning & No dataset-specific training, text-guided, supports creative design & Limited fine control, lower geometry accuracy, CLIP-dependent \\
\hline
11 & DreamHuman (Hu et al., 2023) & Diffusion-based human modeling & Multi-view diffusion, NeRF integration & High realism, strong generalization, integrates multiple AI models & Very high compute cost, slow inference, complex implementation \\
\hline
\end{longtable}
\end{center}

\section{Research Gap Identification}

Based on the comprehensive literature survey, the following research gaps have been identified:

\begin{itemize}
    \item \textbf{Integrated 2D-3D Workflow:} Most existing systems focus either on 2D virtual try-on (FASHN-VTON, VITON, CP-VTON) or 3D avatar generation (PIFuHD, AG3D, PIXIE), but rarely provide both modalities in an integrated workflow accessible via standard web browsers.

    \item \textbf{Accessible Deployment:} State-of-the-art methods like HumanNeRF, DreamHuman, and MetaHuman require significant computational resources and specialized software ecosystems, limiting their deployment in consumer-facing web applications.

    \item \textbf{Single-Image Reconstruction with Detail Preservation:} While methods like PIFuHD and PIXIE work with single images, they often struggle with complete texture preservation or require parametric constraints that limit clothing fidelity.

    \item \textbf{Parallel Processing Architecture:} Existing systems typically process 2D and 3D modalities sequentially, forcing users to wait for complete pipeline execution before receiving any feedback.

    \item \textbf{Real-time Web Rendering:} The challenge of streaming complex 3D meshes to browsers and rendering them interactively with WebGL remains underexplored in academic literature.

    \item \textbf{Texture Mapping Integration:} No existing system effectively projects 2D garment textures from virtual try-on onto reconstructed 3D meshes for photorealistic visualization.
\end{itemize}

STYLMORPH addresses these gaps by combining FASHN-VTON for high-fidelity 2D try-on with PIFuHD for detailed 3D reconstruction, implementing a parallel pipeline architecture that minimizes perceived wait time, and delivering both modalities through a scalable MERN stack architecture with Three.js-based WebGL rendering.

\chapter{PROPOSED SYSTEM}

In the proposed system, STYLMORPH employs a bifurcated parallel processing pipeline designed to balance speed and depth of visualization. Upon uploading a full-body image, the system simultaneously dispatches data to two distinct pipelines. The primary pipeline ports the image to a remote Google Colaboratory environment via a secure tunneling connection, where the FASHN-VTON framework performs garment warping and synthesis to produce a photorealistic 2D try-on result. The secondary pipeline runs PIFuHD locally within a Python subprocess to generate a dense 3D 'obj' mesh encoding the user's full body topology.

\section{Modules}

STYLMORPH is architected as four core modules, each leveraging modern web technologies and machine learning integration to deliver a seamless virtual try-on experience. These modules work in concert to transform a single 2D photograph into a personalized 3D avatar with realistic garment visualization.

\subsection{User Interface Module}

This module provides an intuitive web-based interface developed using React.js for user interaction. It handles user registration, login, image upload, product selection, and outfit curation. The interface displays generated 3D models via Three.js and allows real-time manipulation of the virtual try-on experience. Cloudinary ensures rapid delivery of 3D assets, while responsive design ensures compatibility across devices and platforms.

\subsection{Body Measurement Processing}

While the primary input is a single 2D photograph processed by PIFuHD, this module processes and validates optional anthropometric measurements provided by users to enhance database accuracy. It includes validation rules and measurement normalization, storing structured numerical inputs (Height, Chest, Waist, etc.) in MongoDB.

\subsection{Machine Learning Module}

This core module executes the two-stage processing pipeline. For 2D try-on, the FASHN-VTON inference framework runs on an external Colab GPU environment due to local infrastructure constraints, handling high-fidelity 2D garment transfer. For 3D reconstruction, the PIFuHD model runs locally via its serialized checkpoint (pifuhd.pt), performing monocular implicit 3D reconstruction and outputting .obj mesh files.

\subsection{3D Rendering and Asset Management}

Generated meshes are automatically uploaded to Cloudinary, with secure asset URLs linked to user Avatar documents in MongoDB. The frontend uses @react-three/fiber and Three.js to provide real-time, hardware-accelerated 360-degree WebGL rendering of 3D avatars directly in the browser without any plugins.

\section{Gantt Chart}

Gantt charts are useful for planning and scheduling projects. They help you assess how long a project should take, determine the resources needed, and plan the order in which you'll complete tasks. They're also helpful for managing the dependencies between tasks.

\begin{figure}[h]
    \centering
    \includegraphics[width=1.2\textwidth]{gantt chart.png}
    \caption{gantt chart}
    \label{fig:Image}
\end{figure}

\section{Cost Estimation}

The purpose of cost estimation is to predict the quantity, cost, and price of the resources required to complete a job within the project scope. Our project falls under the categories of organic.

Variables of organic are:
\[a_b = 2.4\]
\[b_b = 1.05\]
\[c_b = 2.5\]
\[d_b = 0.38\]

\begin{itemize}
    \item EFFORT, \(E = a_b(KLOC)^{b_b} = 2.4 \times 4.5^{1.05} = 11.36\) PM
    \item DURATION, \(D = c_b(E)^{d_b} = 2.5 \times 11.36^{0.38} = 6.2\) Months
    \item STAFFING \(=\) Effort / Duration \(= 11.36 / 6.2 = 1.83\) Persons
    \item Total Cost (at Rs. 40,000 per person-month) \(= 40,000 \times 1.83 \times 6.2 = Rs. 453,840\)
\end{itemize}

\section{Advantages}

\begin{itemize}
    \item \textbf{Two-Stage Workflow:} Combines instant 2D visualization with detailed 3D avatar generation
    \item \textbf{High-Fidelity Avatars:} Generates detailed 3D meshes with clothing wrinkles and hairstyles using PIFuHD
    \item \textbf{Single Image Input:} Requires only one 2D photograph, eliminating need for complex multi-camera setups
    \item \textbf{Parallel Pipeline Architecture:} Ensures both 2D and 3D processing run concurrently, minimizing perceived wait time
    \item \textbf{Cloud-Based Delivery:} Cloudinary CDN enables rapid streaming of 3D assets
    \item \textbf{Browser-Based Rendering:} Three.js WebGL eliminates need for plugins or downloads
    \item \textbf{Reduced Return Rates:} Better purchase decisions through accurate visualization
    \item \textbf{Enhanced Customer Confidence:} Increased trust in online purchases
    \item \textbf{Real-time Processing:} Instant visualization and interactive manipulation
    \item \textbf{Scalable Architecture:} Microservices design capable of handling multiple users simultaneously
\end{itemize}

\chapter{SOFTWARE REQUIREMENT SPECIFICATION}

\subsection{Purpose}
The purpose of STYLMORPH is to develop an advanced AI-powered virtual try-on system that revolutionizes online fashion retail. The system enables customers to visualize how clothing items would look on their specific body type by providing both 2D garment previews and personalized 3D avatars from single 2D photographs, thereby reducing return rates and enhancing shopping satisfaction.

\subsection{Product Scope}
STYLMORPH is designed for integration with e-commerce platforms, fashion retailers, and individual consumers seeking enhanced online shopping experiences. The system serves both B2B and B2C markets, providing virtual try-on capabilities that bridge the gap between physical and digital shopping through state-of-the-art monocular 3D reconstruction and 2D garment transfer.

\section{Overall Description}

\subsection{Product Perspective}
STYLMORPH represents a standalone virtual try-on platform built on the MERN stack with specialized Python microservices for machine learning processing. The system processes user inputs including single 2D photographs and optional body measurements to generate personalized 3D avatars using PIFuHD and 2D try-on results using FASHN-VTON. Generated meshes are delivered via Cloudinary and rendered in the browser using Three.js. By providing accurate virtual try-on experiences, STYLMORPH enhances customer engagement and reduces product returns for retailers.

\subsection{Product Functions}
\begin{itemize}
    \item \textbf{2D Virtual Try-On:} High-fidelity garment transfer using FASHN-VTON preserving textures and details
    \item \textbf{Personalized 3D Avatar Creation:} Generates high-fidelity body models from single 2D images using PIFuHD
    \item \textbf{Parallel Pipeline Execution:} Concurrent 2D and 3D processing for optimal user experience
    \item \textbf{Real-time Garment Mapping:} Simulates clothing items on personalized avatars using Three.js
    \item \textbf{Cloud-Based Asset Delivery:} Cloudinary CDN for rapid 3D mesh streaming
    \item \textbf{Digital Wardrobe Management:} Allows users to curate and save outfit combinations
\end{itemize}

\subsection{User Classes and Characteristics}
\begin{itemize}
    \item \textbf{End Consumers:} Individual shoppers seeking better online fitting experiences
    \item \textbf{Fashion Retailers:} E-commerce businesses aiming to reduce returns and enhance engagement
    \item \textbf{Fashion Brands:} Companies looking to showcase products through virtual try-on
    \item \textbf{System Administrators:} Technical staff managing the platform and user data
\end{itemize}

\subsection{Operating Environment}
\begin{itemize}
    \item \textbf{Hardware:} Standard computing devices with modern browsers supporting WebGL
    \item \textbf{Software:} Web-based platform compatible with major browsers, backend using Node.js/Express.js and Python microservices
    \item \textbf{Cloud Infrastructure:} Cloudinary for asset delivery, MongoDB Atlas for database, Google Colab for FASHN-VTON execution

\end{itemize}

\subsection{Design and Implementation Constraints}
\begin{itemize}
    \item \textbf{Privacy Compliance:} Must adhere to data protection regulations (GDPR, CCPA)
    \item \textbf{Performance Requirements:} ML processing offloaded to Python child processes and external Colab environment
    \item \textbf{Scalability:} Support for multiple concurrent users with Cloudinary CDN
    \item \textbf{Browser Compatibility:} Support for major web browsers with WebGL capabilities
\end{itemize}

\subsection{Assumptions and Dependencies}
\begin{itemize}
    \item \textbf{User Cooperation:} Users will upload clear, full-body photographs
    \item \textbf{Internet Connectivity:} Stable internet connection required for optimal performance
    \item \textbf{Hardware Capabilities:} Adequate device performance for WebGL 3D rendering
    \item \textbf{Data Availability:} Access to pre-trained PIFuHD checkpoints and FASHN-VTON framework
\end{itemize}

\section{External Interface Requirements}

\subsection{User Interfaces}
\begin{itemize}
    \item \textbf{Web Application:} Responsive React.js interface for desktop and mobile users
    \item \textbf{Image Upload:} Interface for uploading single 2D photographs
    \item \textbf{3D Viewer:} Interactive Three.js model visualization with rotation and zoom capabilities
    \item \textbf{2D Try-On Display:} Visualization of FASHN-VTON synthesized results
    \item \textbf{Digital Wardrobe:} Browsing interface for available 3D clothing items
\end{itemize}

\subsection{Hardware Interfaces}
\begin{itemize}
    \item \textbf{Camera Support:} For capturing user photographs
    \item \textbf{Processing Unit:} Adequate CPU/GPU for WebGL 3D rendering
\end{itemize}

\subsection{Software Interfaces}
\begin{itemize}
    \item \textbf{E-commerce APIs:} Integration with product catalogs
    \item \textbf{Cloud Services:} Cloudinary for asset storage and delivery
    \item \textbf{AI Frameworks:} PyTorch for PIFuHD, Google Colab for FASHN-VTON
\end{itemize}

\subsection{Communications Interfaces}
\begin{itemize}
    \item \textbf{REST APIs:} For system integration and data exchange
    \item \textbf{HTTPS:} Secure data transmission for user privacy
    \item \textbf{Secure Tunneling:} For communication with Google Colab environment
\end{itemize}

\section{System Features}

\subsection{Description and Priority}
This is the highest prioritized feature. It is the main step in the application.

\subsection{Stimulus/Response Sequences}
Normal Flow of Events:
\begin{itemize}
    \item The user uploads a full-body photograph
    \item The system processes the image through parallel pipelines
    \item 2D try-on results and 3D avatar are generated
    \item The matched videos and 3D models are displayed on screen
\end{itemize}

\subsection{Functional Requirements}
\begin{itemize}
    \item \textbf{REQ1:} The system shall allow users to create and manage personal profiles
    \item \textbf{REQ2:} The system shall accept and process single 2D photograph uploads
    \item \textbf{REQ3:} The system shall generate 2D virtual try-on results using FASHN-VTON
    \item \textbf{REQ4:} The system shall generate personalized 3D avatars using PIFuHD
    \item \textbf{REQ5:} The system shall upload generated meshes to Cloudinary and store URLs in MongoDB
    \item \textbf{REQ6:} The system shall render 3D avatars in browser using Three.js
    \item \textbf{REQ7:} The system shall provide interactive 360-degree avatar visualization
    \item \textbf{REQ8:} The system shall support digital wardrobe management and outfit curation
\end{itemize}

\section{Other Nonfunctional Requirements}

\subsection{Performance Requirements}
This app requires laptops with basic processor and browser Chrome with latest version. The system must support multiple concurrent users with efficient ML processing offloaded to Python child processes and external Colab execution. 3D rendering should be smooth and interactive via Three.js WebGL.

\subsection{Safety Requirements}
The web app will be regularly updated for fixing bugs and adding new features. In case of problems, the user can contact the app support team for any service.

\subsection{Security Requirements}
User data must be encrypted both in transit and at rest. The system must comply with data protection regulations and implement proper access controls. This web app requires microphone access for voice capture.

\subsection{Software Quality Attributes}
This web app for virtual try-on provides users with both simple and advanced features. Due to its well-designed and easy-to-use interface, it can be used by both experts and typical users. The application mainly focuses on ordinary people with basic knowledge of using smartphones or any browsing devices.

\chapter{SYSTEM DESIGN}

\section{Architecture Diagram}

An architecture diagram represents the structural design of a system, showing how different modules interact with each other. It provides a high-level overview of the system components, data flow, and processing pipelines involved in the application.

\begin{figure}[h]
    \centering
    \includegraphics[width=0.8\textwidth]{Architecture_Diagram(fixed) (1).png}
    \caption{System Architecture Diagram}
    \label{fig:architecture_flow}
\end{figure}

The architecture diagram shown in Fig.~\ref{fig:architecture_flow} illustrates the complete workflow of the STYLMORPH system. The process begins with the user interacting with the application through the UI module, where they provide login/signup details, upload their image, and enter a wearable item link.

The user data and uploaded image are stored in the database, while the wearable link is sent to the web scraper module, which extracts the garment image from the given URL. These inputs are then passed through the processing layer via API communication.

The system executes two parallel pipelines within the 3D Generator module. The first pipeline uses the FASHN-VTON algorithm to generate realistic 2D virtual try-on images. The second pipeline uses the PIFuHD algorithm to reconstruct a detailed 3D human model from the uploaded image.

Finally, the generated outputs are delivered back to the frontend. The user can visualize the results as high-quality 2D try-on images and an interactive 3D avatar model.

\subsection{Module-wise Explanation}

\subsubsection{1. User Module}
This module represents the end-user who interacts with the system. The user provides necessary inputs such as authentication details, personal images, and garment selection links.

\subsubsection{2. UI Module}
The UI module handles all frontend interactions. It includes:
\begin{itemize}
    \item Login/Signup interface
    \item Image upload functionality
    \item Wearable item link input
\end{itemize}
It communicates with the backend to send user inputs and display outputs.

\subsubsection{3. Database Module}
The database stores user information, uploaded images, garment data, and processed outputs. It ensures persistent storage and efficient retrieval of data.

\subsubsection{4. Web Scraper Module}
This module extracts garment images from the provided wearable item link. It enables integration with external e-commerce platforms by dynamically fetching product images.

\subsubsection{5. API / Processing Layer}
This layer manages communication between frontend and backend components. It handles requests, coordinates data flow, and triggers processing pipelines.

\subsubsection{6. 3D Generator Module}
This is the core processing unit of the system, consisting of:
\begin{itemize}
    \item \textbf{FASHN-VTON (2D Virtual Try-On):} Generates realistic images of the user wearing selected garments.
    \item \textbf{PIFuHD (3D Reconstruction):} Converts the 2D user image into a detailed 3D avatar model.
\end{itemize}

\subsubsection{7. Output Module}
The output module presents the final results to the user. It includes:
\begin{itemize}
    \item 2D try-on images
    \item Interactive 3D avatar model
\end{itemize}
This allows users to visualize garment fitting in both 2D and 3D formats.

\section{Flowchart}

A flowchart is a type of diagram that represents a workflow or process. A flowchart can also be defined as a diagrammatic representation of an algorithm, a step-by-step approach to solving a task.

\begin{figure}[h]
    \centering
    \includegraphics[width=0.5\textwidth]{Activity diagram (1).png}
    \caption{Activity Diagram}
    \label{fig:Activity}
\end{figure}


The flowchart shows the steps as boxes of various kinds, and their order by connecting the boxes with arrows. Fig 7.1, showing the flowchart of our system, explains the complete working of our system. User can enter the web app, then upload image data. The image is processed through parallel pipelines - one for 2D FASHN-VTON try-on and another for PIFuHD 3D reconstruction. The output from both pipelines is stored in Cloudinary and displayed to the user.

\section{Use Case Diagram}

\begin{figure}[h]
    \centering
    \includegraphics[width=0.5\textwidth]{Usecase.png}
    \caption{Usecase Diagram}
    \label{fig:Usecase}
\end{figure}

The use case diagram illustrates the core interactions for a user within the Stylmorph application. The primary actor, the User, engages with the system to create a personalized virtual avatar. The process begins with the user logging in or signing up to access the platform. Subsequently, the user can upload their own image, allowing the system to automatically fetch their skin tone. To build the virtual wardrobe, the user can also upload images of wearable items (such as clothes or accessories). Finally, the user can save the configured avatar along with their style preferences, which the system then stores for future sessions. This diagram effectively maps the complete user journey from authentication to the creation and storage of a personalized style profile.

\section{Sequence Diagram}

\begin{figure}[h]
    \centering
    \includegraphics[width=1\textwidth]{SEQUENCE (1).png}
    \caption{Sequence Diagram}
    \label{fig:Sequence}
\end{figure}

A Sequence Diagram is an interaction diagram that details how operations are carried out, what messages are sent and when. Sequence diagrams are organized according to time. The time progresses as you go down the page. The objects involved in the operation are listed from left to right according to when they take part in the message sequence.

In Fig 7.3, we can see that the user authenticates and uploads a photograph to the Node.js backend, then it delegates to Python modules for FASHN-VTON on Colab and PIFuHD locally, and finally corresponding 3D meshes and 2D images are fetched from Cloudinary and shown as output.

\section{Data Flow Diagram}

DFD graphically representing the functions, or processes, which capture, manipulate, store, and distribute data between a system and its environment and between components of a system. The visual representation makes it a good communication tool between User and System designer. Structure of DFD allows starting from a broad overview and expand it to a hierarchy of detailed diagrams.

\subsection{LEVEL 0 - Context Diagram}

\begin{figure}[h]
    \centering
    \includegraphics[width=0.5\textwidth]{dfdlevel0 (3).png}
    \caption{Level-0 DFD Diagram}
    \label{fig:Level-0 DFD}
\end{figure}

In Fig 7.4, the user uploads a photograph as input to the STYLMORPH system to produce 2D try-on results and 3D avatar visualization understandable by the user.

\subsection{LEVEL 1 - Major Processes}

\begin{figure}[h]
    \centering
    \includegraphics[width=0.5\textwidth]{dfdlevel1 (2).png}
    \caption{Level-1 DFD Diagram}
    \label{fig:Level-1 DFD}
\end{figure}

The Level 1 DFD decomposes the system into major functional processes including User Authentication, 2D Try-On Processing (FASHN-VTON on Colab), 3D Reconstruction (PIFuHD), Asset Management, and 3D Rendering. This will be seen in Fig 7.5.

\subsection{LEVEL 2 - Detailed Processes}

\begin{figure}[h]
    \centering
    \includegraphics[width=0.7\textwidth]{dfdlevel2(original) (3).png}
    \caption{Level-2 DFD Diagram}
    \label{fig:Level-2 DFD}
\end{figure}

The Level 2 DFD provides detailed decomposition of the two-stage processing pipeline, showing how the FASHN-VTON pipeline on Google Colab processes input images for 2D try-on results, while the PIFuHD pipeline generates 3D mesh files, followed by Cloudinary upload and URL storage in MongoDB, as in Fig 7.6.

\section{Entity-Relationship Diagram}

The Entity-Relationship Diagram defines the database schema for STYLMORPH, showing entities, attributes, and relationships.

\begin{table}[H]
\centering
\caption{Avatar Data Schema Structure}
\label{tab:avatar-schema}
\begin{tabular}{|l|p{8cm}|}
\hline
\textbf{Field} & \textbf{Description} \\
\hline
UserId & Unique identifier linking the avatar to a user account \\
ImageUrl & Cloudinary URL of the user's uploaded 2D source photo \\
VtonResults & Array of 2D try-on images from FASHN-VTON with garment metadata \\
PifuhdUrl & Cloudinary URL of the generated PIFuHD .obj file \\
Wearables & Array of digital clothing items (Name, URL, Thumbnail) applied to the avatar \\
SavedSets & Categorized configurations of wearables saved as full outfits \\
Measurements & Optional structured numerical inputs (Height, Chest, Waist, etc.) \\
\hline
\end{tabular}
\end{table}

\begin{figure}[h]
    \centering
    \includegraphics[width=0.8\textwidth]{er.jpg}
    \caption{ER Diagram}
    \label{fig:ER}
\end{figure}

\chapter{IMPLEMENTATION}

\section{Dataset}
The STYLMORPH framework leverages a combination of widely-recognized benchmarks and user-contributed datasets to ensure robust performance across various body types and garment styles. For the 2D Virtual Try-On module, we reference datasets like VITON-HD, while the 3D reconstruction pipeline utilizes pre-trained PIFuHD checkpoints. The application dynamically processes user-uploaded high-resolution images, which are stored in Cloudinary's cloud infrastructure. This real-time dataset allows users to see personalized results based on their own physical attributes, moving beyond static libraries.

\subsection{Algorithms Used}

\subsubsection{PIFuHD Algorithm (3D Model Generation)}
The PIFuHD (Pixel-Aligned Implicit Function for High-Resolution) algorithm is used to generate a detailed 3D human model from a single 2D image. It learns an implicit representation of the human body to achieve high-resolution reconstruction.

Initially, a Convolutional Neural Network (CNN) extracts important features from the input image, such as body shape, clothing details, and textures. For every point in the 3D space, the algorithm predicts whether the point lies inside or outside the human body using an implicit function.

To maintain spatial consistency, each 3D point is projected back onto the 2D image plane, ensuring alignment with the original image pixels. After evaluating all points, the Marching Cubes algorithm is applied to convert the implicit representation into a continuous 3D surface.

The final output is a high-resolution \texttt{.obj} mesh that captures fine details such as wrinkles, hair, and body structure.



\subsubsection{FASHN-VTON Algorithm (2D Virtual Try-On)}
The FASHN-VTON algorithm is used to generate realistic 2D virtual try-on results by digitally applying garments onto a user’s image. It is based on a diffusion model, specifically a Multimodal Diffusion Transformer (MMDiT), which progressively transforms noise into a realistic dressed image.

To ensure proper garment fitting, Thin Plate Spline (TPS) transformation is used to warp the clothing image according to the user’s body shape and pose. Additionally, pose estimation techniques such as DWPose detect key body landmarks (e.g., shoulders, arms, and waist), enabling accurate alignment of garments.

\section{Implementation Details}
The system architecture follows a modern full-stack approach. The frontend is built using React to provide a declarative and responsive user interface. The backend utilizes Node.js and Express, integrated with MongoDB for persistent storage and Cloudinary for media management.

\subsection{User and Avatar Modules}
We decouple basic authentication from user-specific physical profiles for better scalability. The \texttt{User} model handles registration and login, while the \texttt{Avatar} model stores complex measurement data and asset URLs for 3D avatars and try-on results.

\begin{verbatim}
// Avatar.js - Schema for 3D Models and Measurements
const mongoose = require("mongoose");

const AvatarSchema = new mongoose.Schema({
  userId: { type: mongoose.Schema.Types.ObjectId, ref: "User", required: true },
  measurements: {
    height: { type: Number },
    chest: { type: Number },
    waist: { type: Number },
    hips: { type: Number },
    shoulder: { type: Number },
    source: { type: String, enum: ['auto', 'manual'] }
  },
  pifuhdUrl: { type: String }, // OBJ file URL
  imageUrl: { type: String },  // Processed person image
  wearables: [{
    url: String,
    name: String,
    thumbnail: String
  }],
  tryonResults: [{
    url: String,
    category: String,
    createdAt: { type: Date, default: Date.now }
  }]
});

module.exports = mongoose.model("Avatar", AvatarSchema);
\end{verbatim}

\subsection{2D Virtual Try-On Module}
The 2D Virtual Try-On (VTON) module manages asynchronous communication with a specialized image processing server via an ngrok tunnel. It downloads the user avatar and garment images, executes the processing on a GPU-enabled infrastructure, and stores the resulting try-on image in the cloud.

\begin{verbatim}
// avatar.js - Controller for 2D VTON
exports.viewWearable = async (req, res) => {
    const avatar = await Avatar.findOne({ userId: req.user.id });
    const wearable = avatar.wearables.find(w => w._id.toString() === wearableId);
    
    // Convert images to Base64 for processing server
    const personBase64 = await downloadAsBase64(avatar.imageUrl);
    const garmentBase64 = await downloadAsBase64(wearable.url);

    // Call external GPU endpoint (VTON-HD/FASHN based)
    const tryonResponse = await fetch("https://vton-api-tunnel.ngrok-free.dev/tryon", {
      method: "POST",
      body: JSON.stringify({
        person_image: personBase64,
        garment_image: garmentBase64,
        category: "tops"
      }),
    });

    const tryonData = await tryonResponse.json();
    const resultUrl = await uploadToCloudinary(tryonData.result_image);
    
    avatar.tryonResults.push({ url: resultUrl, category: "tops" });
    await avatar.save();
};
\end{verbatim}

\subsection{3D Model Generation (PIFuHD)}
This module orchestrates the generation of a 3D avatar from a single 2D image using PIFuHD. Once the high-density OBJ model is generated, a custom Python script extracts precise body measurements based on the user's real-world height.

\begin{verbatim}
// avatar.js - 3D Reconstruction and Measurement Extraction
exports.generateModels = async (req, res) => {
    const userHeight = parseFloat(req.body.height);
    const pifuOutputPath = `./temp/pifuhd_${Date.now()}.obj`;

    // 1. Run PIFuHD Reconstruction
    await runPythonScript("run_pifuhd.py", imageUrl, pifuOutputPath);

    // 2. Extract Measurements from OBJ
    const measJson = await runPythonScript("extract_measurements.py", 
                                           pifuOutputPath, userHeight);
    const extractedMeasurements = JSON.parse(measJson);

    // 3. Update Profile
    const avatar = await Avatar.findOneAndUpdate(
        { userId: req.user.id },
        { 
          pifuhdUrl: await uploadSmplObj(pifuOutputPath),
          measurements: { ...extractedMeasurements, source: "auto" }
        }
    );
};
\end{verbatim}

\subsection{Garment Mapping and Simulation (Three.js)}
The simulation module uses Three.js within the React environment to render the reconstructed 3D avatar. It implements sophisticated lighting systems (Warm/Cool Rim lights) and an OrbitControls interface to allow users to interact with their digital twin.

\begin{verbatim}
// RealisticAvatar3D.jsx - Three.js Implementation
const RealisticAvatar3D = ({ modelUrl, measurements }) => {
    useEffect(() => {
        const scene = new THREE.Scene();
        const camera = new THREE.PerspectiveCamera(55, w / h, 0.01, 200);
        const renderer = new THREE.WebGLRenderer({ antialias: true });
        
        // Advanced Skin Lighting
        scene.add(new THREE.AmbientLight(0xffeedd, 0.60));
        const keyLight = new THREE.DirectionalLight(0xfff5e0, 1.4);
        keyLight.position.set(-3, 9, 6);
        scene.add(keyLight);

        // Load PIFuHD OBJ Model
        const loader = new OBJLoader();
        loader.load(modelUrl, (object) => {
            object.traverse((child) => {
                if (child.isMesh) {
                    child.material = new THREE.MeshStandardMaterial({
                        color: 0xffffff, roughness: 0.7
                    });
                }
            });
            scene.add(object);
        });
    }, [modelUrl]);
};
\end{verbatim}


\chapter{SOFTWARE TESTING}

Software testing is a process used for verifying the correctness, completeness and quality of the developed software. Software is built out of sub-systems that are composed of modules, which in turn are composed of procedures and functions. The sequence of testing activities performed for the STYLMORPH system is as below:

\begin{itemize}
    \item Unit Testing
    \item Integration Testing
    \item System Testing
    \item Functional Testing
\end{itemize}

\section{Unit Testing}

There exist a number of components in every sub-system. Every component is tested using respective test procedures. Each component is tested individually based on their needs. Unit test focuses verification effort on the smallest unit of the software design component. Each module can be tested using the following two Strategies:

\begin{enumerate}
    \item \textbf{Black Box Testing:} In this strategy some test cases are generated as input conditions that fully execute all functional requirements for the program.
    \item \textbf{White Box Testing:} In this the test cases are generated on the logic of each module by drawing flow graphs of that module and logical decisions are tested on all the cases.
\end{enumerate}

\begin{table}[H]
\centering
\caption{Unit Testing}
\label{tab:unit-testing}
\begin{tabular}{|p{1.5cm}|p{1.5cm}|p{1.5cm}|p{1.5cm}|p{1.5cm}|p{1.5cm}|}
\hline
\textbf{Test Case ID} & \textbf{Test Objective} & \textbf{Test Steps} & \textbf{Test Data} & \textbf{Expected Result} & \textbf{Pass/Fail} \\
\hline
M101 & To ensure that user can upload image & Open website and upload image & User photograph & Image upload should complete successfully & Pass \\
\hline
M102 & To ensure 2D try-on generation & Upload image with garment & Test images & 2D try-on result should be generated & Pass \\
\hline
M103 & To ensure 3D avatar generation & Upload image and process & Test photograph & 3D mesh should be created & Pass \\
\hline
M104 & To ensure video display & Process multiple words & Sentence input & Single video should be displayed & Pass \\
\hline
\end{tabular}
\end{table}

\section{Integration Testing}

Integration testing is the process of testing the interface between two software units or modules. It's focus on determining the correctness of the interface. The purpose of the integration testing is to expose faults in the interaction between integrated units. Once all the modules have been unit tested, integration testing is performed.

\begin{table}[H]
\centering
\caption{Integration Testing}
\label{tab:integration-testing}
\begin{tabular}{|p{1.5cm}|p{3cm}|p{3cm}|p{2.5cm}|p{1.5cm}|}
\hline
\textbf{Test Case ID} & \textbf{Test Objective} & \textbf{Test Steps} & \textbf{Expected Outcome} & \textbf{Pass/Fail} \\
\hline
P01 & To ensure all modules communicate & Run all modules and observe output & All modules communicate correctly, proper output shown & Pass \\
\hline
\end{tabular}
\end{table}

\section{System Testing}

System testing is the stage of implementation. Testing is the process of executing a program with the intent of finding an error. A good test case is one that has a high probability of finding a yet undiscovered error.

System testing of software or hardware is testing conducted on a complete, integrated system to evaluate the system's compliance with its specified requirements. System testing falls within the scope of black-box testing, and as such, should require no knowledge of the inner design of the code or logic.

\section{Functional Testing}

Involves in-house testing of the entire system before delivery to the user. Its aim is to satisfy the user the system meets all requirements of the client's specifications. Functional testing means testing the application against business requirements. Functional testing is executed using the functional specifications given by the client or by the design specifications according to use cases given by the design team. Role of functional testing is to validating the behavior of an application.

\section{Test Case}

Test cases define how to test a system. A test case is a singular set of actions for a tester to perform that validates a specific aspect of a product or application functionality. If the test fails, the result might be a software defect that the organization can triage.

\begin{table}[H]
\centering
\caption{Test Case}
\label{tab:test-case}
\begin{tabular}{|p{1.5cm}|p{2.5cm}|p{1.5cm}|p{1.5cm}|p{1.5cm}|p{1.5cm}|}
\hline
\textbf{Test Case ID} & \textbf{Test Scenario} & \textbf{Test Steps} & \textbf{Test Data} & \textbf{Expected Data} & \textbf{Pass/Fail} \\
\hline
M101 & Upload single image & Upload photograph through interface & Test image JPG & Successful upload & Pass \\
\hline
M102 & Generate 2D try-on & Select garment and process & Person + garment images & 2D try-on result displayed & Pass \\
\hline
M103 & Generate 3D avatar & Process uploaded image & Test photograph & 3D mesh created and displayed & Pass \\
\hline
\end{tabular}
\end{table}

\section{Test Result}

\subsection{Test Case 1: Image Upload}

Fig 9.1 shows the result of Test Case 1. Any image that is given as input is successfully uploaded and processed. So test case 1 is successfully passed.

\begin{figure}[h]
    \centering
    \includegraphics[width=0.7\textwidth]{Image Upload.jpeg}
    \caption{Image Upload}
    \label{fig:Image}
\end{figure}

\subsection{Test Case 2: 2D Try-On Generation}

The outcome of Test Case 2 is shown in Fig 9.2. Any garment and person image entered as input is correctly processed and 2D try-on result is displayed. As a consequence, test case 2 was performed successfully.

\begin{figure}[h]
    \centering
    \includegraphics[width=0.7\textwidth]{2D Try-On Generation.jpeg}
    \caption{2D Try-On Result}
    \label{fig:Image}
\end{figure}

\subsection{Test Case 3: 3D Avatar Generation}

The outcome of Test Case 3 is shown in Figure 9.3. Any photograph entered as input is correctly processed and 3D avatar is generated. As a result, test case 3 has been performed successfully.


\begin{figure}[h]
    \centering
    \includegraphics[width=0.7\textwidth]{ 3D Avatar Generation.jpeg}
    \caption{ 3D Avatar Generation}
    \label{fig:Image}
\end{figure}




\chapter{RESULT AND ANALYSIS}

\section{Result}

\begin{figure}[h]
    \centering
    \includegraphics[width=0.7\textwidth]{Home Page .jpeg}
    \caption{ Home Page}
    \label{fig:Image}
\end{figure}

We built Frontend, which is a react application. React is a JavaScript package that allows you to build interactive web applications. Declarative, quick, and adaptable user interfaces. You may use it to create complex structures. Components are small, self-contained chunks of code that make up user interfaces. As shown in Fig. 10.1, there is a button that says "Upload Image". It assists in the uploading of photographs for processing.

\begin{figure}[h]
    \centering
    \includegraphics[width=0.7\textwidth]{ 2D Try-On Generation.jpeg}
    \caption{ 2D Try-On Generation}
    \label{fig:Image}
\end{figure}

The parallel pipeline architecture successfully decouples the slow 2D Colab inference from the local 3D reconstruction, ensuring both run concurrently. The FASHN-VTON pipeline running over Google Colab produced high-fidelity try-on images preserving fine garment textures, logos, and folds across varied body poses, as shown in Fig 10.2.

\begin{figure}[h]
    \centering
    \includegraphics[width=0.7\textwidth]{3D Avatar Generation.jpeg}
    \caption{ 3D Avatar Output}
    \label{fig:Image}
\end{figure}

The PIFuHD pipeline effectively extracted detailed 3D body topology from single monocular images, capturing clothing wrinkles and hair details in the output mesh. Three.js rendered the resulting meshes smoothly in the browser at stable frame rates, even for complex meshes exceeding 50,000 polygon faces, as shown in Fig 10.3.


\begin{figure}[h]
    \centering
    \includegraphics[width=0.7\textwidth]{ Screenshot of Result.jpeg}
    \caption{ Screenshot of Result}
    \label{fig:Image}
\end{figure}
The parallel pipeline architecture successfully decouples the slow 2D Colab inference from the local 3D reconstruction, ensuring both run concurrently. The FASHN-VTON pipeline running over Google Colab produced high-fidelity try-on images preserving fine garment textures, logos, and folds across varied body poses. The PIFuHD pipeline effectively extracted detailed 3D body topology from single monocular images, capturing clothing wrinkles and hair details in the output mesh. Three.js rendered the resulting meshes smoothly in the browser at stable frame rates.

\section{Performance Analysis}

Table 10.1 presents the measured execution latencies across all pipeline stages during system testing.

\begin{table}[H]
\centering
\caption{Pipeline Execution Latency and Output Metrics}
\label{tab:latency}
\begin{tabular}{|p{2.5cm}|p{2.5cm}|p{2.5cm}|p{2.5cm}|}
\hline
\textbf{Pipeline Stage} & \textbf{Latency} & \textbf{Output} & \textbf{User Impact} \\
\hline
FASHN-VTON (Colab) & 7-8 min & JPEG Image & 2D garment visualization \\
PIFuHD 3D Reconstruction & 3-4 min & .obj Mesh & Full 360° avatar fit view \\
Three.js WebGL Rendering & < 2 sec & Canvas (WebGL) & Interactive 3D exploration \\
\hline
\end{tabular}
\end{table}

The 2D FASHN-VTON pipeline requires 7-8 minutes owing to the remote Colab execution overhead and network transfer latency. The 3D PIFuHD reconstruction completes locally in 3-4 minutes. Both pipelines run concurrently, so the user can inspect the 3D avatar while the 2D try-on is still processing, minimizing perceived wait time.

\chapter{CONCLUSION AND SCOPE OF FUTURE WORK}

In a variety of situations, a sign language interpreter is quite helpful. Anyone may use this technology to communicate in sign language at schools, colleges, hospitals, universities, airlines, courts, and anyplace else. It facilitates communication between those who are normally hearing and others who are hard of hearing. Deaf people may participate in all of the activities that hearing people can, from daily interactions to information access, thanks to this technology. Our proposed model will successfully convert the given input audio into video frames.

STYLMORPH successfully demonstrates the feasibility of integrating state-of-the-art 2D virtual try-on (FASHN-VTON) with explicit monocular 3D reconstruction (PIFuHD) within a scalable, consumer-accessible web application. The MERN-stack architecture, augmented by Three.js WebGL rendering and Cloudinary CDN delivery, proves that high-fidelity fashion visualization is viable for real-world e-commerce deployment. By providing both immediate 2D garment previews and detailed 3D avatars, Stylmorph offers a significant advancement over existing single-modality approaches.

Many improvements along this route can be made as and when the ISL Dictionary grows. The words in the ISL are small, so many improvements can be made by adding new words to their dictionary to increase their breadth. In addition, text-and-speech integration can be done on a project to enable better communication techniques that will allow users to convert Text into Indian Sign Language by hand-input text. The future work lays in the direction of improving the current model with a custom speech recognition system than Google API.

According to the need definition, we've completed the project. The modules have been carefully planned and thoroughly researched to meet the project's needs. The report has been completed with a lot of effort and with our guide's full assistance and advice.

\section{Future Directions}

Future work will focus on three key areas:

\begin{itemize}
    \item \textbf{Migrating FASHN-VTON Pipeline:} Move the FASHN-VTON pipeline to a dedicated GPU inference server to reduce latency from minutes to seconds
    \item \textbf{Texture Mapping Integration:} Project FASHN-VTON 2D texture maps directly onto PIFuHD 3D mesh surfaces to produce fully textured photorealistic avatars
    \item \textbf{Physics-Based Cloth Simulation:} Implement algorithms within the Three.js ecosystem to calculate collision layers between generated human meshes and digital garments, ensuring realistic fabric drape and avoiding geometric clipping
    \item \textbf{Real-time Computer Vision Pre-processor:} Integrate client-side photo quality validation to guide users in taking optimal setup photos
    \item \textbf{Augmented Reality Integration:} Enable real-time virtual try-on using smartphone cameras
    \item \textbf{Sustainability Features:} Show environmental impact of fashion choices by visualizing reduced return rates and carbon footprint
\end{itemize}

The success of STYLMORPH demonstrates the growing importance of virtual try-on technology in the future of retail. As consumer expectations evolve and technology advances, systems like STYLMORPH will play an increasingly crucial role in shaping how people discover, try, and purchase fashion items. The foundation established by this project provides a solid platform for continued innovation in personalized digital fashion experiences.

\begin{thebibliography}{99}

\bibitem{1} Saito, S., Huang, Z., Natsume, R., et al. "PIFuHD: Multi-Level Pixel-Aligned Implicit Function for High-Resolution 3D Human Digitization." IEEE/CVF Conference on Computer Vision and Pattern Recognition (CVPR), 2020.

\bibitem{2} Loper, M., Mahmood, N., Romero, J., Pons-Moll, G., \& Black, M. J. "SMPL: A Skinned Multi-Person Linear Model." ACM Transactions on Graphics (TOG), 34(6), 1-16, 2015.

\bibitem{3} Dong, Z., et al. "AG3D: Learning to Generate 3D Avatars from 2D Images." IEEE Transactions on Pattern Analysis and Machine Intelligence, 45(8), 11234-11248, 2023.

\bibitem{4} Ma, Q., et al. "CAPE: Clothed Auto-Personalized Avatars." ACM Transactions on Graphics, 41(4), 1-15, 2022.

\bibitem{5} Batool, R., \& Mou, J. "Virtual Try-On Technology: A Comprehensive Review and Research Agenda." International Journal of Human-Computer Interaction, 39(5), 1023-1045, 2023.

\bibitem{6} Hosen, M. S., et al. "Mastering 3D Modeling in Blender." International Journal of Computer Graphics, 15(2), 45-62, 2023.

\bibitem{7} Jaleel, M. U. "Impact of 3D Product Displays on E-Commerce Consumer Behavior." International Journal of E-Commerce Studies, 14(3), 201-220, 2023.

\bibitem{8} Yang, H., Yao, H., \& Zhou, S. "Accurate 3D Human Body Measurement from Monocular Images Using Silhouette Analysis." Pattern Recognition Letters, 142, 18-25, 2021.

\bibitem{9} Kanazawa, A., Black, M. J., Jacobs, D. W., \& Malik, J. "End-to-End Recovery of Human Shape and Pose." IEEE/CVF Conference on Computer Vision and Pattern Recognition (CVPR), 7122-7131, 2018.

\bibitem{10} Pavlakos, G., et al. "Learning to Estimate 3D Human Pose and Shape from a Single Color Image." IEEE/CVF Conference on Computer Vision and Pattern Recognition (CVPR), 459-468, 2018.

\bibitem{11} Bogo, F., et al. "Keep It SMPL: Automatic Estimation of 3D Human Pose and Shape." European Conference on Computer Vision (ECCV), 561-578, 2016.

\bibitem{12} Lassner, C., Pons-Moll, G., \& Gehler, P. "A Generative Model of People in Clothing." International Conference on Computer Vision (ICCV), 853-862, 2017.

\bibitem{13} Yang, Z., et al. "Graphonomy: Universal Human Parsing via Graph Transfer Learning." IEEE/CVF Conference on Computer Vision and Pattern Recognition (CVPR), 7450-7459, 2020.

\bibitem{14} Pons-Moll, G., et al. "Dyna: A Model of Dynamic Human Shape in Motion." ACM Transactions on Graphics (TOG), 34(4), 120:1-120:14, 2015.

\bibitem{15} Guan, P., et al. "Drape: Dressing Any Person." ACM Transactions on Graphics (TOG), 31(4), 35:1-35:10, 2012.

\bibitem{16} FASHN Research. "FASHN-VTON: High-Fidelity Virtual Try-On with Texture Preservation." FASHN Research Technical Report, 2024. Available: https://fashn.ai/research/vton

\bibitem{17} Han, X., et al. "VITON: An Image-based Virtual Try-on Network." IEEE/CVF Conference on Computer Vision and Pattern Recognition (CVPR), 7543-7552, 2018.

\bibitem{18} Wang, B., et al. "Toward Characteristic-Preserving Image-based Virtual Try-On Network." European Conference on Computer Vision (ECCV), 589-604, 2018.

\bibitem{19} Software Engineering (3rd edition): Book by Aggarwal K.K. Singh Yogesh




\end{thebibliography}

\end{document}